\documentclass[a4paper,11pt]{article}

\usepackage[finnish]{babel}
\usepackage{xcolor}
%\usepackage[framemethod=TikZ]{mdframed}
\usepackage[utf8]{inputenc}
\usepackage[margin=2cm]{geometry}
\usepackage{amsfonts,amsmath,amssymb,amsthm,enumitem}
\usepackage{microtype}
\usepackage{pgf}
\usepackage{tikz}
\usepackage{environ}
\usetikzlibrary{arrows,automata}

\setenumerate{listparindent=\parindent}

\newtheorem*{claim}{Väite}
\newtheorem*{lemma}{Lemma}
\theoremstyle{definition}
\newtheorem*{definition}{Määritelmä}

\newtheorem{exercisethm}{Tehtävä}
%\colorlet{shadecolor}{cyan!8}
%\newenvironment{exercise}
%{\begin{mdframed}[roundcorner=10pt]\begin{exer}}
        %{\end{exer}\end{mdframed}}

\newcommand{\set}[1]{{\left\{ #1 \right\}}}
\newcommand{\ceil}[1]{{\left\lceil#1\right\rceil}}
\newcommand{\Nat}{\mathbb{N}}

\tikzstyle{mybox} = [draw=none, fill=cyan!9,
    rectangle, rounded corners, inner sep=10pt, inner ysep=20pt]

\NewEnviron{exercise}%
{  \begin{tikzpicture}
        \node [mybox] (box) {
            \begin{minipage}{0.9\textwidth}
                \begin{exercisethm}
                    \BODY
                \end{exercisethm}
            \end{minipage}
        };
    \end{tikzpicture}
}


\title{Pumppauslemmaopas}
\author{Juhana Laurinharju \and Jani Rahkola}

\begin{document}

\maketitle


\section*{Pumppauslemma}

Jokaisella säännöllisellä kiellellä on seuraava \emph{pumppauslemmana} tunnettu
ominaisuus.

\begin{definition}[pumppauslemma]
    Olkoon $A$ säännöllinen kieli. Tällöin $A$:lla on jokin pumppauspituus $p
    \in \Nat$, $p > 0$. Nyt kaikki tätä pidemmät merkkijonot $s \in A$, $|s|
    \geq p$ voidaan jakaa kolmeen osaan $s = xyz$ siten, että seuraavat kolme
    ehtoa ovat voimassa.

    \begin{enumerate}
        \item $xy^iz \in A$ jokaisella $i \in \Nat$, erityisesti myös kun $i = 0$.
        \item $|y| > 0$, eli toistettava osa $y$ ei saa olla tyhjä merkkijono $\varepsilon$.
        \item $|xy| \leq p$
    \end{enumerate}
\end{definition}

Otetaan tästä konkreettinen esimerkki. Kieli $A = L(a^*b^*)$ on säännöllinen ja
sillä on täten jokin pumppauspituus $p$. Tämän kielen kohdalla eräs mahdollinen
pumppauspituus on $p = 2$. Tarkastellaan jotain riittävän pitkää kielen $A$
merkkijonoa. Valitaan $s = abb$. Nyt pumppauslemman nojalla löytyy \emph{jokin}
jako $s = xyz$, jolle yllä olevat kolme ehtoa pätevät.

Tarkastellaan merkkijonon $s$ mahdollisia jakoja.

\begin{itemize}
    \item Kokeillaan ensin seuraavaa jakoa:
        \begin{align*}
             x               & = \varepsilon,  \\
             y               & = ab            \\
            \textrm{ ja }  z & = b
        \end{align*}
        %
        Nyt ehdot 2 ja 3 ovat voimassa, mutta ensimmäinen ehto ei täyty, sillä
        esimerkiksi merkkijono $xyyz = ababb$ ei kuulu kieleen $A$.

    \item Ensimmäinen jako ei siis täyttänyt kaikkia pumppauslemman ehtoja.
        Pumppauslemma ei kuitenkaan takaa, että nämä ehdot täyttyisivät
        jokaisella jaolla. Ainoa tae on se, että löytyy jokin ehdot täyttävä
        jako. Tällainen on esimerkiksi seuraava jako:
        %
        \begin{align*}
            x              & = a \textrm{,} \\
            y              & = b \\
            \textrm{ja } z & = b
        \end{align*}
        Nyt $|xy| = 2 \leq p$, $|y| = 1 > 0$. Entä miltä näyttää merkkijono
        $xy^iz$? Tarkastellaan tätä ensin $i$:n arvoilla $0$, $1$ ja $2$:
        %
        \begin{align*}
            xy^0z & = xz  = ab \in A \\
            xy^1z & = xyz = abb \in A \\
            xy^2z & = xyyz = abbb \in A
        \end{align*}
        %
        Ja yleisessä tapauksessa, kun $i \in \Nat$, niin
        %
        \begin{equation*}
            xy^iz = ab^ib = ab^{i+1} \in L(a^*b^*) = A
        \end{equation*}
\end{itemize}

\begin{exercise}
    Olkoon
    %
    \begin{equation*}
        A = L((ab)^*
    \end{equation*}
    %
    säännöllinen kieli. Tällä kielellä on pumppauspituus $p = 2$. Valitaan
    kielestä merkkijono
    %
    \begin{equation*}
        s = ababab
    \end{equation*}
    %
    joka on pidempi kuin pumppauspituus $p = 2$.  Anna kaikki merkkijonon $s$
    jaot $s = xyz$ jotka täyttävät pumppauslemman ehdot $2$ ja $3$, eli 
    %
    \begin{align*}
        |y|  & > 0 \textrm{ ja} \\
        |xy| & \leq p \text{.}
    \end{align*}
\end{exercise}

\begin{exercise}
    Mitkä edellisen tehtävän jaoista $s = xyz$ toteuttavat pumppauslemman
    ensimmäisen ehdon? Ensimmäinen ehto on
    %
    \begin{equation*}
        xy^iz \in A
    \end{equation*}
    %
    eli keskikohtaa $y$ voi toistaa.
\end{exercise}

\subsection*{Pumppauslemman käyttäminen}

Pumppauslemma on hyödyllinen koska sillä voidaan näyttää \emph{joitakin} kieliä
epäsäännöllisiksi. Eräs tällainen kieli on jo tuttu ja epäsäännölliseksi
osoitettu $A = \set{0^n1^n \mid n \in \Nat}$  Osoitetaan tämän kielen
epäsäännöllisyys pumppauslemman avulla.
%
\begin{claim}
    Kieli $A = \set{0^n1^n \mid n \in \Nat}$ on epäsäännöllinen.
\end{claim}
\begin{proof}
    asd
\end{proof}

\end{document}
